\documentclass[11pt,letterpaper]{article}
\usepackage[margin=1in,headheight=15pt]{geometry}
\usepackage[T1]{fontenc}
\usepackage{newtxtext}
\usepackage{mathtools,amsthm}
\usepackage{newtxmath}
\usepackage{microtype}
\usepackage[dvipsnames]{xcolor}
\usepackage{booktabs,array,longtable}
\usepackage{enumitem}
\usepackage{fancyhdr}
\usepackage{titlesec}
\usepackage[most]{tcolorbox}
\usepackage{listings}
\usepackage{hyperref}
\definecolor{DeepBlue}{HTML}{17324D}
\definecolor{Teal}{HTML}{176B70}
\definecolor{PaleBlue}{HTML}{F0F5F8}
\definecolor{QuietGray}{HTML}{52606D}
\hypersetup{colorlinks=true,linkcolor=DeepBlue,citecolor=Teal,urlcolor=Teal,
 pdftitle={Arithmetic of a Self-Iterating Generating Function},
 pdfauthor={Research report prepared with ChatGPT},bookmarksnumbered=true}
\titleformat{\section}{\large\bfseries\color{DeepBlue}}{\thesection}{0.7em}{}
\titleformat{\subsection}{\normalsize\bfseries\color{Teal}}{\thesubsection}{0.7em}{}
\pagestyle{fancy}
\fancyhf{}
\fancyhead[L]{\small\color{QuietGray}Arithmetic of a self-iterating generating function}
\fancyhead[R]{\small\color{QuietGray}A396807}
\fancyfoot[C]{\small\thepage}
\renewcommand{\headrulewidth}{0.3pt}
\setlength{\parskip}{4pt}
\setlength{\parindent}{1.1em}
\setlist{nosep,leftmargin=1.5em}
\newtheorem{theorem}{Theorem}[section]
\newtheorem{lemma}[theorem]{Lemma}
\newtheorem{proposition}[theorem]{Proposition}
\newtheorem{corollary}[theorem]{Corollary}
\theoremstyle{definition}
\newtheorem{definition}[theorem]{Definition}
\theoremstyle{remark}
\newtheorem{remark}[theorem]{Remark}
\newcommand{\Z}{\mathbb Z}
\newcommand{\Q}{\mathbb Q}
\newcommand{\F}{\mathbb F}
\newcommand{\iter}[2]{#1^{\langle #2\rangle}}
\newcommand{\coeff}[1]{[x^{#1}]}
\newcommand{\M}{\mathcal M}
\DeclareMathOperator{\Tr}{Tr}
\DeclareMathOperator{\lcm}{lcm}
\lstset{basicstyle=\small\ttfamily,columns=fullflexible,breaklines=true,
 frame=single,rulecolor=\color{PaleBlue},showstringspaces=false,
 xleftmargin=0.3em,xrightmargin=0.3em}
\title{\vspace{-1.0cm}\color{DeepBlue}\bfseries
 Arithmetic of a Self-Iterating\\[3pt] Generating Function\\[7pt]
 \large Proofs of the A396807 congruences, rational prime-power lifts,\\
 and an exact period of 46,860 for the last two digits}
\author{\normalsize Research report prepared with ChatGPT}
\date{September 19, 2026}
\begin{document}
\maketitle
\thispagestyle{fancy}

\begin{abstract}
Let $A=x+O(x^2)$ be the formal power series defined by
$A=x+\iter{A}{5}\iter{A}{6}$, where angle brackets denote functional
iteration.  We prove the two congruences recorded as conjectures in OEIS
A396807: $\coeff{n}A\equiv1\pmod {10}$ and
$\coeff{n}\iter{A}{k}\equiv k^{n-1}\pmod {10}$ for every integer iterate
$k$ and every $n\geq1$.  The key is uniqueness over arbitrary coefficient
rings and the substitution $\M(x)=x/(1-x)$.

The main extension is an arithmetic theory of the entire family
$F=x+\iter{F}{r}\iter{F}{s}$.  Its largest uniform M\"obius-congruence
modulus is exactly $\gcd(r+s-1,rs)$.  For each prime dividing this modulus,
we explicitly invert the linearized functional equation using a cyclic
trace.  Consequently, every prime-power reduction is rational and its
coefficient sequence is eventually periodic.  For A396807 we calculate
the reductions modulo $4$ and $25$ and prove that the last two decimal
digits have minimal period $46860$, starting at $n=1$.  We also determine
exact compositional orders, construct $2$-adic and $5$-adic iterates,
prove zero complex radius of convergence, and settle closely related
congruences in two neighboring OEIS entries.  All proofs are accompanied
by exact, reproducible computation; numerical tests are not used as a
substitute for the all-degree arguments.
\end{abstract}

\begin{tcolorbox}[colback=PaleBlue,colframe=DeepBlue,boxrule=0.5pt,
 arc=1pt,left=8pt,right=8pt,top=6pt,bottom=6pt]
\textbf{Status and scope.}
The OEIS records retrieved on September 19, 2026 still label the target
statements as conjectures.  This report supplies self-contained proofs
and further deductions.  A limited search did not locate a separate
proof of the target statements; this is not a claim of exhaustive
literature coverage or priority.  The report has not been externally
peer reviewed or formalized in a proof assistant.
\end{tcolorbox}

\clearpage
\tableofcontents
\newpage

\section{The problem, the answer, and the notation}\label{sec:problem}

\subsection{A concrete conjecture}
Paul D. Hanna's OEIS entry A396807, contributed on June 16, 2026,
defines the sequence $a(n)$ by the ordinary formal generating function
\begin{equation}\label{eq:target}
 A(x)=\sum_{n\geq1}a(n)x^n
      =x+\iter{A}{5}(x)\iter{A}{6}(x).
\end{equation}
The entry explicitly uses the exponents for \emph{iteration}, not
ordinary multiplication, and records the two congruences addressed
below \cite{oeis807,oeis807internal}.

Throughout the report,
\[
 \iter{F}{0}(x)=x,\qquad
 \iter{F}{j+1}=F\circ\iter{F}{j}.
\]
For negative $j$, $\iter{F}{j}$ is an iterate of the compositional
inverse.  Such an inverse exists over any commutative ring for a series
$F=x+O(x^2)$.  Ordinary powers, for example $F(x)^2$, retain their usual
multiplicative meaning.  This distinction is essential: replacing the
right side of \eqref{eq:target} by $x+A(x)^{11}$ gives an entirely
different problem.

The first coefficients, also independently reproduced by the supplied
program, are
\[
 1,\ 1,\ 11,\ 201,\ 4721,\ 129671,\ 3976201,\ 132554541,
 \ 4724742051,\ 178052854981,\ldots.
\]
All congruences between power series in this report are coefficientwise.
They do not assert convergence at a numerical value of $x$.

\begin{theorem}[The two target congruences]\label{thm:target}
The unique solution of \eqref{eq:target} in $x+x^2\Z[[x]]$ satisfies
\begin{equation}\label{eq:maincong}
 A(x)\equiv\frac{x}{1-x}\pmod {10},\qquad
 \iter{A}{k}(x)\equiv\frac{x}{1-kx}\pmod {10}
 \quad(k\in\Z).
\end{equation}
In particular, for $n\geq1$,
\[
 a(n)\equiv1\pmod {10},\qquad
 \coeff{n}\iter{A}{k}\equiv k^{n-1}\pmod {10}.
\]
At $k=0,n=1$, the last display means that the coefficient is $1$;
equivalently, use the convention $0^0=1$ in this one coefficient formula.
The modulus $10$ is the largest positive integer for which all the
coefficients of $A$ equal $1$ modulo that integer.
\end{theorem}

The proof occupies Sections \ref{sec:formal}--\ref{sec:modulus} and does
not require computation.  The stronger results are summarized here to
make their logical status clear.

\begin{center}
\begin{tabular}{@{}p{0.25\textwidth}p{0.57\textwidth}p{0.08\textwidth}@{}}
\toprule
Result & Conclusion & Section\\
\midrule
Exact modulus & A necessary and sufficient condition for the whole
 two-iterate family, including an integer weight & \ref{sec:modulus}\\[3pt]
Rational lifting & Rational reduction modulo every $p^e$ at each
 prime selected by the exact modulus & \ref{sec:lifting}\\[3pt]
Decimal digits & Explicit formulas modulo $4$, $20$, $25$, and $100$;
 minimal two-digit period $46860$ & \ref{sec:period}\\[3pt]
Iteration & Exact order $p^e$ modulo $p^e$ and well-defined
 $p$-adic-time iterates & \ref{sec:orders}\\[3pt]
Related entries & Modulo-$6$ results for A396797; all modulo-$8$
 iterate formulas for A396798, with a corrected starting index
 & \ref{sec:neighbors}\\[3pt]
Growth & A factorial lower bound along each parity and zero
 complex radius of convergence & \ref{sec:growth}\\
\bottomrule
\end{tabular}
\end{center}

\section{Formal existence, uniqueness, and reduction}\label{sec:formal}

\subsection{A coefficientwise contraction}
Let $R$ be a commutative ring with identity.  Congruence modulo $x^N$
means equality of the coefficients in degrees less than $N$.  The
$x$-adic topology simply packages this coefficientwise stabilization.
For fixed nonnegative integers $r,s$ and $c\in R$, set
\begin{equation}\label{eq:T}
 T_c(F)=x+c\,\iter{F}{r}\iter{F}{s},\qquad F\in x+x^2R[[x]].
\end{equation}

\begin{lemma}[One extra coefficient]\label{lem:contraction}
If $F\equiv G\pmod {x^N}$ with $N\geq2$, then
\[
 \iter{F}{j}\equiv\iter{G}{j}\pmod {x^N}\quad(j\geq0),
 \qquad T_c(F)\equiv T_c(G)\pmod {x^{N+1}}.
\]
\end{lemma}
\begin{proof}
Composition by a series of order one preserves divisibility by $x^N$.
Also, if $U-V$ is divisible by $x^N$ and $U(0)=V(0)=0$, then
$H(U)-H(V)$ is divisible by $x^N$: apply the identity that
$U^m-V^m$ is divisible by $U-V$ to each coefficient of $H$.
Induction on $j$, using
\[
 F\circ\iter{F}{j}-G\circ\iter{G}{j}
 =(F-G)\circ\iter{F}{j}
   +G\circ\iter{F}{j}-G\circ\iter{G}{j},
\]
proves the assertion about iterates.  The difference of the two products
in \eqref{eq:T} is
\[
 (\iter{F}{r}-\iter{G}{r})\iter{F}{s}
 +\iter{G}{r}(\iter{F}{s}-\iter{G}{s}).
\]
Every iterate on the right has zero constant term, including the
zeroth iterate $x$.  Each summand is therefore divisible by $x^{N+1}$.
No division, cancellation of ring elements, or field hypothesis is used.
\end{proof}

\begin{theorem}[Well-posedness over arbitrary rings]\label{thm:unique}
There is a unique $F\in x+x^2R[[x]]$ satisfying
\begin{equation}\label{eq:generalweighted}
 F=x+c\,\iter{F}{r}\iter{F}{s}.
\end{equation}
Its coefficient of $x^2$ is $c$.  The construction commutes with every
homomorphism of coefficient rings.
\end{theorem}
\begin{proof}
Starting at $F_0=x$, put $F_{j+1}=T_c(F_j)$.  The first difference is
$c x^2$, and Lemma \ref{lem:contraction} shows inductively that
$F_{j+1}-F_j$ is divisible by $x^{j+2}$.  Every coefficient stabilizes,
defining a series $F$.  Each fixed-degree coefficient of a composition
or product depends on only finitely many coefficients of its inputs.
Passing to the stabilized coefficients gives $F=T_c(F)$.

Any two solutions agree modulo $x^2$.  Repeated application of the lemma
makes them agree modulo every power of $x$, hence they are equal.
The product in \eqref{eq:generalweighted} begins with $x^2$, proving the
claim about its coefficient.  Ring homomorphisms preserve all the
finite sums and products used to construct a fixed coefficient, so
construction and coefficient reduction commute.
\end{proof}

\subsection{Why negative iterates cause no difficulty}
If $F=x+\sum_{n\geq2}u_nx^n$, the coefficient of $x^n$ in
$F(G)$, where $G=x+\sum_{n\geq2}v_nx^n$, is $v_n$ plus a polynomial
with integer coefficients in previously determined $v_j$ and the
$u_j$.  Setting $F(G)=x$ therefore determines each $v_n$ uniquely without
dividing by a nonunit.  The same argument constructs a left inverse;
associativity of composition makes the left and right inverses equal.

It follows that inversion, positive iteration, and negative iteration
all commute with reduction modulo an integer.  In particular, the
uniqueness theorem may be used directly over $\Z/10\Z$, although this
ring is not a field.  This is the exact logical step that promotes a
rational substitution from an observation to a proof for every
coefficient.

\section{The M\"obius substitution and the exact modulus}\label{sec:modulus}

\subsection{A calculation valid in every ring}
For $c$ in a coefficient ring define
\[
 M_c(x)=\frac{x}{1-cx}.
\]
The denominator has constant term $1$, so the expression is a formal
power series over the same ring.  Direct composition gives
\begin{equation}\label{eq:mobiusgroup}
 M_u\circ M_v=M_{u+v},\qquad
 \iter{M_c}{j}(x)=\frac{x}{1-jcx}\quad(j\in\Z).
\end{equation}
Write $S=r+s$ and $P=rs$.  Clearing only denominators that are formal
units gives the identity
\begin{equation}\label{eq:residual}
 M_c-T_c(M_c)
 =\frac{-c^2(S-1)x^3+c^3P x^4}
 {(1-cx)(1-rcx)(1-scx)}.
\end{equation}
Thus the two numbers $c^2(S-1)$ and $c^3P$ govern the entire
coefficientwise congruence.

\begin{theorem}[Sharp uniform modulus]\label{thm:sharp}
Let $r,s\geq0$ and $c\in\Z$.  Let
$F_{r,s,c}=x+\sum_{n\geq2}u_nx^n$ be the unique integral solution of
\eqref{eq:generalweighted}.  For every positive integer $m$,
\begin{equation}\label{eq:iff}
 F_{r,s,c}\equiv M_c\pmod m
 \quad\Longleftrightarrow\quad
 m\mid c^2(r+s-1)\ \text{ and }\ m\mid c^3rs.
\end{equation}
Equivalently,
\begin{equation}\label{eq:gcdideal}
 \gcd_{n\geq2}\bigl|u_n-c^{n-1}\bigr|
 =\gcd\bigl(c^2(r+s-1),c^3rs\bigr).
\end{equation}
We use $\gcd(0,0)=0$; a zero right side means exact equality, and hence
congruence for every modulus.  For $c=1$ the sharp modulus is
\begin{equation}\label{eq:D}
 D_{r,s}=\gcd(r+s-1,rs).
\end{equation}
\end{theorem}
\begin{proof}
If the two divisibilities hold, \eqref{eq:residual} shows that $M_c$ is
a solution of the defining equation over $\Z/m\Z$.  By
Theorem \ref{thm:unique}, it is the reduction of $F_{r,s,c}$.

For necessity, the first four coefficients suffice.  The elementary
iteration expansion is
\[
 \iter{F}{j}=x+jc x^2+
   \bigl(j u_3+j(j-1)c^2\bigr)x^3+O(x^4).
\]
Substitution in the defining equation gives
\begin{equation}\label{eq:firstfour}
 F_{r,s,c}=x+c x^2+c^2Sx^3+
          c^3(2S^2-S-P)x^4+O(x^5).
\end{equation}
Consequently, its discrepancies from $M_c$ in degrees $3$ and $4$ are
\[
 c^2(S-1),\qquad
 c^3\bigl((S-1)(2S+1)-P\bigr).
\]
If both vanish modulo $m$, subtracting an integer multiple of the first
from the second yields $m\mid c^3P$.  These same two discrepancies have
gcd $\gcd(c^2(S-1),c^3P)$.  Since that gcd divides every discrepancy by
the sufficient direction, \eqref{eq:gcdideal} follows.
\end{proof}

\begin{corollary}[All integer iterates]\label{cor:iterates}
Whenever \eqref{eq:iff} holds,
\[
 \iter{F_{r,s,c}}{k}(x)\equiv\frac{x}{1-kcx}\pmod m
 \qquad(k\in\Z).
\]
For $n\geq1$, the corresponding coefficient is $(kc)^{n-1}$, with
the degree-one convention stated in Theorem \ref{thm:target}.
\end{corollary}
\begin{proof}
Reduction commutes with iteration and inversion, and
\eqref{eq:mobiusgroup} gives the iterates of the reduced series.
\end{proof}

\subsection{The target conjecture and a useful subfamily}
For $(r,s,c)=(5,6,1)$,
\[
 D_{5,6}=\gcd(10,30)=10.
\]
More concretely,
\[
 (1-5x)(1-6x)=1-11x+30x^2\equiv1-x\pmod {10},
\]
so that
\[
 x+\frac{x}{1-5x}\frac{x}{1-6x}
 \equiv x+\frac{x^2}{1-x}=\frac{x}{1-x}\pmod {10}.
\]
Uniqueness proves the first target congruence, and
Corollary \ref{cor:iterates} proves the second.  Finally $a(3)-1=10$,
so a larger common modulus is impossible.  This proves
Theorem \ref{thm:target}.

For adjacent iterate indices $(r,s)=(q,q+1)$ with $q\geq1$,
\begin{equation}\label{eq:adjacentD}
 D_{q,q+1}=\gcd(2q,q(q+1))
 =\begin{cases}2q,&q\text{ odd},\\q,&q\text{ even}.\end{cases}
\end{equation}
The modulus is therefore not uniformly $2q$: the parity of $q$ matters.
This explains simultaneously the moduli $6$, $4$, and $10$ in the
three examples studied here.

\begin{remark}[Scaling]
If $F_{r,s,1}=x+\sum_{n\geq2}a_nx^n$, then
\[
 F_{r,s,c}(x)=x+\sum_{n\geq2}c^{n-1}a_nx^n.
\]
For nonzero $c$, this is the identity
$F_{r,s,c}(x)=c^{-1}F_{r,s,1}(cx)$, interpreted coefficientwise;
it also holds at $c=0$ without any division.  Direct substitution and
uniqueness prove the assertion.
\end{remark}

\section{A cyclic-trace inverse for the linearized equation}\label{sec:trace}

The uniform congruence is elementary.  To go beyond it, we must control
how the self-iteration equation changes under a small perturbation.
The crucial fact is that the reduced M\"obius map has a finite orbit.
This section gives an explicit inverse, rather than invoking an
unspecified lifting theorem.

\subsection{The cyclic substitution and its trace}
Fix a prime $p\mid D_{r,s}$.  In $\F_p$, the conditions
$r+s=1$ and $rs=0$ force the residues of $r,s$ to be $0,1$ in some
order.  Interchanging them does not change the functional equation.
We may therefore write
\begin{equation}\label{eq:alphabeta}
 r=p\alpha,\qquad s=1+p\beta,
 \qquad \alpha,\beta\in\Z_{\geq0}.
\end{equation}
In formulas over $\F_p$, $\alpha$ and $\beta$ mean their residues.
Put
\[
 M_j(x)=\frac{x}{1-jx},\qquad M=M_1,\qquad M_0=x.
\]
Equation \eqref{eq:mobiusgroup} implies $M_p=x$ over $\F_p$.
For a formal series or a rational function regular at zero define
\begin{equation}\label{eq:trace}
 \Tr_p f=\sum_{j=0}^{p-1}f(M_j(x)).
\end{equation}
The trace is invariant under $x\mapsto M(x)$, by cyclically reindexing
the sum.  If $B$ is invariant, then
$\Tr_p(Bf)=B\Tr_p(f)$.

\begin{lemma}[Trace of the coordinate]\label{lem:tracex}
Over $\F_p(x)$,
\begin{equation}\label{eq:tracex}
 \Tr_p x=-\frac{x^p}{1-x^{p-1}}.
\end{equation}
Consequently, with
\begin{equation}\label{eq:C}
 C(x)=\alpha x+\beta M(x),
\end{equation}
we have
\begin{equation}\label{eq:traceC}
 \Tr_p C=-\frac{(\alpha+\beta)x^p}{1-x^{p-1}}.
\end{equation}
\end{lemma}
\begin{proof}
In $\F_p[t]$,
$\prod_{j\in\F_p}(t-j)=t^p-t$.  Taking its logarithmic derivative gives
\[
 \sum_{j\in\F_p}\frac1{t-j}=-\frac1{t^p-t}.
\]
Substituting $t=1/x$ proves \eqref{eq:tracex}.  The trace of $M(x)$
equals that of $x$, again by cyclic reindexing, proving
\eqref{eq:traceC}.
\end{proof}

\subsection{First variation of a functional iterate}
Work temporarily over $\F_p[\varepsilon]/(\varepsilon^2)$.  Write a
perturbation in the normalized form
\[
 F_\varepsilon=M+\varepsilon M^2w.
\]
The factor $M^2$ simplifies the chain rule.

\begin{lemma}[Iteration derivative]\label{lem:variation}
For every $k\geq0$,
\begin{equation}\label{eq:variation}
 \iter{F_\varepsilon}{k}
 =M_k+\varepsilon M_k^2\sum_{j=0}^{k-1}w(M_j(x)).
\end{equation}
The sum is empty at $k=0$.
\end{lemma}
\begin{proof}
The formula is immediate at $k=0,1$.  In passing from $k$ to $k+1$, the
perturbation already present in the inner iterate is multiplied by
$M'(M_k)$.  Since
\[
 M'(M_k)M_k^2=M_{k+1}^2,
\]
and the new outer perturbation contributes
$M_{k+1}^2w(M_k)$, the partial sum acquires exactly one further term.
\end{proof}

Let $T(F)=x+\iter{F}{r}\iter{F}{s}$ and denote by
\[
 Lh=h-(DT)_M h
\]
the linearization of $F-T(F)$ at $M$.  In the notation
\eqref{eq:alphabeta}, the sums in \eqref{eq:variation} are
\[
 \sum_{j=0}^{r-1}w(M_j)=\alpha\Tr_p w,
 \qquad
 \sum_{j=0}^{s-1}w(M_j)=\beta\Tr_p w+w.
\]
Using $M_r=x$ and $M_s=M$ in $\F_p[[x]]$ gives the central identity
\begin{equation}\label{eq:L}
 L(M^2w)=M^2(1-x)\bigl(w-C\Tr_p w\bigr).
\end{equation}
For clarity, the two product-derivative terms, divided by $M^2$, are
$x(1-x)\alpha\Tr_p w$ and $x(\beta\Tr_p w+w)$; subtracting them from
$w$ yields \eqref{eq:L}.

\begin{theorem}[Explicit inverse]\label{thm:inverseL}
For any $R\in x^3\F_p[[x]]$, set
\begin{equation}\label{eq:inverseL}
 d=\frac{R}{M^2(1-x)},\qquad
 w=d+C\frac{\Tr_p d}{1-\Tr_p C},\qquad h=M^2w.
\end{equation}
Then $h\in x^3\F_p[[x]]$ is the unique solution of $Lh=R$.
If $R$ is rational and regular at zero, then so is $h$.
\end{theorem}
\begin{proof}
The equation $Lh=R$ is equivalent to $w-C\Tr_p w=d$.
Taking the trace and using the invariance of $\Tr_p w$ gives
\[
 (1-\Tr_p C)\Tr_p w=\Tr_p d.
\]
The factor $1-\Tr_p C$ has constant term $1$, so it is invertible.
Substituting the resulting trace into the original equation gives
\eqref{eq:inverseL}, and proves uniqueness as well.

Since $R$ starts in degree at least $3$, $d$ starts in degree at least
$1$.  Both $C$ and $\Tr_p d$ have zero constant term.  Thus $w$ starts
in degree at least $1$ and $M^2w$ in degree at least $3$.
The operations in \eqref{eq:inverseL} are finite rational operations
and substitution by the rational functions $M_j$, all regular at zero.
They preserve rationality and regularity.
\end{proof}

\section{Rational reduction modulo every selected prime power}\label{sec:lifting}

\subsection{The lifting theorem}
Here a series over $\Z/p^e\Z$ is called rational when it can be written
as $P(x)/Q(x)$ with polynomial numerator and denominator and with
$Q(0)$ a unit.  This definition is meaningful even though the coefficient
ring is not a field.

\begin{theorem}[Rational prime-power lifts]\label{thm:rationallifts}
Let $F=F_{r,s,1}$ and let $p$ be a prime dividing $D_{r,s}$.
For every $e\geq1$, the coefficientwise reduction of $F$ modulo $p^e$
is rational.  The proof gives a finite recursive construction of a
rational representative.
\end{theorem}
\begin{proof}
Let $\Z_{(p)}$ denote the rational numbers whose denominators are not
divisible by $p$.  Work with rational functions having coefficients in
$\Z_{(p)}$ and a denominator whose constant term is a $p$-adic unit.
Their expansions at zero lie in $\Z_{(p)}[[x]]$.

The initial representative is $f_1=M$, by Theorem \ref{thm:sharp}.
Inductively, suppose that the rational function $f_e$ agrees with $F$
modulo $p^e$ and has its coefficients of $x$ and $x^2$ exactly equal to
$1$, with zero constant term.  Its residual is rational and satisfies
\[
 T(f_e)-f_e\equiv0\pmod {p^e}.
\]
The quotient
\begin{equation}\label{eq:Re}
 R_e=\frac{T(f_e)-f_e}{p^e}\pmod p
\end{equation}
is a rational element of $x^3\F_p[[x]]$.  To justify the division
algebraically, write the residual over a common denominator with unit
constant term.  Multiplication by this denominator shows that every
coefficient of its polynomial numerator is divisible by $p^e$ in
$\Z_{(p)}$.  Dividing the numerator by $p^e$ preserves the permitted
class of rational functions.

Solve $Lh_e=R_e$ by Theorem \ref{thm:inverseL}.  Choose a rational lift
$\widetilde h_e\in x^3\Z_{(p)}[[x]]$ of $h_e$, with denominator still
having unit constant term, and put
\begin{equation}\label{eq:liftstep}
 f_{e+1}=f_e+p^e\widetilde h_e.
\end{equation}
Because $2e\geq e+1$, terms involving two perturbations vanish modulo
$p^{e+1}$.  Also $f_e\equiv M\pmod p$, so its derivative operator reduces
to $(DT)_M$.  Hence
\[
 T(f_{e+1})-f_{e+1}
 \equiv p^e\bigl(R_e+(DT)_M h_e-h_e\bigr)
 =0\pmod {p^{e+1}}.
\]
Uniqueness over $\Z/p^{e+1}\Z$ identifies $f_{e+1}$ with the reduction
of $F$.  It is rational and still has exactly the prescribed first two
coefficients.  This completes the induction.
\end{proof}

\begin{corollary}[Eventual periodicity]\label{cor:periodic}
For every such $p$ and every $e\geq1$, the coefficient sequence of $F$
is eventually periodic modulo $p^e$.  In particular, A396807 is rational
modulo $10^e$, and its last $e$ decimal digits are eventually periodic,
for every $e\geq1$.
\end{corollary}
\begin{proof}
Normalize a rational denominator to have constant term $1$.
Its coefficients impose a fixed finite-order forward linear recurrence
after the degree of the numerator.  Over a finite ring there are only
finitely many states consisting of the last fixed number of
coefficients.  A repeated state gives a repeated future, proving eventual
periodicity.

For A396807 use $p=2,5$.  Rational representatives modulo $2^e$ and
$5^e$ can have their denominators lifted with constant term exactly $1$.
A Chinese-remainder linear combination is then a rational representative
modulo $10^e$.  The same finite-state argument applies.
\end{proof}

\begin{remark}[What the theorem does not say]
The construction proves rationality separately at each finite modulus.
It does not assert that one rational function over $\Q$ works at every
modulus, nor that its numerator and denominator degrees stay bounded as
$e$ increases.  It supplies no small uniform bound on the periods at
higher powers.  In fact, Section \ref{sec:growth} proves that the
characteristic-zero generating function is not rational.
\end{remark}

\subsection{A closed formula for the first new digit}
The first lifting step is particularly simple because its normalized
residual is exactly $C$.

\begin{theorem}[Universal first correction]\label{thm:firstcorrection}
With $r=p\alpha$ and $s=1+p\beta$, define in $\F_p(x)$
\begin{equation}\label{eq:Hgeneral}
 H_{p;\alpha,\beta}(x)=
 \frac{x^3\bigl(\alpha(1-x)+\beta\bigr)(1-x^{p-1})}
 {(1-x)^3\bigl(1-x^{p-1}+(\alpha+\beta)x^p\bigr)}.
\end{equation}
Then
\begin{equation}\label{eq:firstlift}
 F_{r,s,1}(x)\equiv M(x)+pH_{p;\alpha,\beta}(x)\pmod {p^2}.
\end{equation}
On the right, any integral rational lift of $H$ with unit denominator
constant may be used.
\end{theorem}
\begin{proof}
Identity \eqref{eq:residual}, with $c=1$ and with its sign reversed,
gives
\[
 \frac{T(M)-M}{p}
 \equiv\frac{x^3\bigl(\alpha+\beta-\alpha x\bigr)}{(1-x)^2}
 \pmod p.
\]
Indeed, $(r+s-1)/p=\alpha+\beta$ and
$rs/p=\alpha(1+p\beta)\equiv\alpha\pmod p$.
After division by $M^2(1-x)$, this residual becomes
\[
 d=\alpha x+\beta\frac{x}{1-x}=C.
\]
Formula \eqref{eq:inverseL} therefore simplifies to
$w=C/(1-\Tr_p C)$.  Multiply by $M^2$ and apply
\eqref{eq:traceC} to obtain \eqref{eq:Hgeneral}.
The lifting proof then gives \eqref{eq:firstlift}.
\end{proof}

This formula also explains why Artin--Schreier polynomials appear in a
question that initially seems to concern only the last decimal digit.
Writing $\lambda=\alpha+\beta$, the reciprocal of the new denominator
factor is
\[
 t^p-t+\lambda.
\]
If $\lambda\ne0$ in $\F_p$, any root $\theta$ satisfies
$\theta^{p^j}=\theta-j\lambda$.  Its Frobenius orbit therefore has length
exactly $p$, proving that this polynomial is irreducible of degree $p$.
This elementary observation will be decisive for the exact period.

\section{Explicit refinements and the two-digit period}\label{sec:period}

\subsection{Reductions modulo 4 and 25}
For A396807 the applicable orientations in \eqref{eq:alphabeta} are
\[
 \begin{array}{c|c|c|c}
 p&r&s&(\alpha,\beta)\bmod p\\\hline
 2&6&5&(1,0)\\
 5&5&6&(1,1).
 \end{array}
\]
Substitution into Theorem \ref{thm:firstcorrection} yields the following
particularly convenient representatives:
\begin{align}
 H_2(x)&=\frac{x^3}{1+x^3},\label{eq:H2}\\
 H_5(x)&=\frac{x^3(2-x)(1-x^4)}
 {(1-x)^3(1-x^4+2x^5)}\nonumber\\
 &=\frac{x^3(2-x)(1+x+x^2+x^3)}
 {(1-x)^2(1-x^4+2x^5)}.\label{eq:H5}
\end{align}
The simplification of $H_2$ is performed over $\F_2$; the displayed
rational function is then an integral lift of that result.

\begin{corollary}[Refinements of the target congruence]\label{cor:refinements}
Coefficientwise,
\begin{align}
 A&\equiv M+2H_2\pmod4,\label{eq:mod4}\\
 A&\equiv M+5H_5\pmod {25},\label{eq:mod25}\\
 A&\equiv M+10\frac{x^3}{1-x^3}\pmod {20},\label{eq:mod20}\\
 A&\equiv M+50H_2+80H_5\pmod {100}.\label{eq:mod100}
\end{align}
In particular,
\begin{equation}\label{eq:mod20coeff}
 a(n)\equiv
 \begin{cases}11\pmod {20},&3\mid n,\\1\pmod {20},&3\nmid n,
 \end{cases}\qquad n\geq1.
\end{equation}
\end{corollary}
\begin{proof}
The first two statements are the two first-correction formulas.
Modulo $4$, $2x^3/(1+x^3)=2x^3/(1-x^3)$, while modulo $5$ the correction
is zero.  The Chinese remainder theorem proves \eqref{eq:mod20}.
For \eqref{eq:mod100}, note that $50\equiv2\pmod4$, $50\equiv0\pmod {25}$,
$80\equiv0\pmod4$, and $80\equiv5\pmod {25}$.
\end{proof}

The first few last-two-digit values are a useful check, but not a proof
of periodicity:
\begin{center}
\begin{tabular}{@{}rrrrrrrrrr@{}}
\toprule
$n$&1&2&3&4&5&6&7&8&9\\
$a(n)\bmod100$&01&01&11&01&21&71&01&41&51\\
\midrule
$n$&10&11&12&13&14&15&16&17&18\\
$a(n)\bmod100$&81&81&31&81&61&71&41&61&11\\
\bottomrule
\end{tabular}
\end{center}

\subsection{A denominator criterion for a minimal period}
\begin{lemma}\label{lem:periodcriterion}
Let $H=N/Q\in\F_p(x)$ be in reduced form, with $Q(0)\ne0$ and
$\deg N\leq\deg Q$.  The coefficients $h_n=\coeff{n}H$, indexed by
$n\geq1$, have period $L$ if and only if $Q$ divides $1-x^L$.
\end{lemma}
\begin{proof}
If $h_{n+L}=h_n$ for $n\geq1$, then
\[
 H=h_0+\frac{\sum_{j=1}^Lh_jx^j}{1-x^L},
\]
so the reduced denominator divides $1-x^L$.
Conversely, if it divides, $(1-x^L)H$ is a polynomial of degree at most
$L$.  Comparing coefficients above degree $L$ gives exactly
$h_{n+L}=h_n$ for $n\geq1$.
\end{proof}

\begin{theorem}[Exact two-digit period]\label{thm:46860}
The sequence $(a(n)\bmod100)_{n\geq1}$ is periodic from its first term,
with minimal period
\begin{equation}\label{eq:46860}
 \boxed{46860.}
\end{equation}
The minimal periods modulo $4$ and modulo $25$ are $3$ and $15620$,
respectively.
\end{theorem}
\begin{proof}
Modulo $4$, formula \eqref{eq:mod4} gives the repeating pattern
$1,1,3$, hence minimal period $3$.

Modulo $25$, the period equals the period of the coefficients of $H_5$
over $\F_5$, since $a(n)=1+5\coeff{n}H_5$ in $\Z/25\Z$ for $n\geq1$.
In \eqref{eq:H5} its reduced denominator is
\begin{equation}\label{eq:Q5}
 Q(x)=(1-x)^2Q_5(x),\qquad Q_5(x)=1-x^4+2x^5.
\end{equation}
To check that it is reduced, first evaluate the numerator at $x=1$:
it is $4\ne0$ in $\F_5$.  Next, the reciprocal polynomial
\[
 q(t)=t^5-t+2
\]
is irreducible.  Indeed, if $q(\theta)=0$, then
$\theta^{5^j}=\theta-2j$, whose first return to $\theta$ occurs at $j=5$.
Thus the Frobenius orbit, and hence the minimal polynomial, has degree
$5$.  The numerator of \eqref{eq:H5} is a product of $x^3$, a linear
factor, and a degree-three factor, none of which is divisible by this
irreducible reciprocal factor.

Let $\theta$ be the class of $t$ in $\F_5[t]/(q)$.
The multiplicative order of $\theta$ divides
$5^5-1=3124=4\cdot11\cdot71$.  Reduction of powers using
$\theta^5=\theta-2$ gives
\begin{align}
 \theta^{1562}&=4,\label{eq:cert1}\\
 \theta^{284}&=4+\theta+2\theta^2+4\theta^4,\label{eq:cert2}\\
 \theta^{44}&=2+2\theta+\theta^2+3\theta^3+3\theta^4.\label{eq:cert3}
\end{align}
None is $1$, because a nonzero polynomial of degree less than $5$
cannot vanish at $\theta$.  The three exponents are $3124/2$, $3124/11$,
and $3124/71$; they exclude every possible proper divisor as the order.
Hence $\theta$, and also its reciprocal, has order exactly $3124$.
The displayed residues are small, independently checkable certificates;
the supplied program obtains them by binary powering and reduction.

It follows that $Q_5$ divides $1-x^L$ precisely when $3124\mid L$.
The factor $(1-x)^2$ divides $1-x^L$ precisely when $5\mid L$:
write $L=5^v u$ with $5\nmid u$ and use
$1-x^L=(1-x^u)^{5^v}$.  By Lemma \ref{lem:periodcriterion}, the minimal
period of $H_5$ is therefore
\[
 \lcm(5,3124)=15620.
\]
The degree condition in that lemma is satisfied: both the numerator
and denominator in the reduced expression for $H_5$ have degree $7$.
Thus there is no initial transient for the coefficients indexed from
$n=1$.

A period modulo $100$ must be a period modulo both $4$ and $25$;
conversely, a common period works by the Chinese remainder theorem.
The minimal period is consequently
$\lcm(3,15620)=46860$.
\end{proof}

\subsection{The general first-lift period mechanism}
The same proof gives more than this numerical example.

\begin{proposition}\label{prop:generalperiod}
Assume $\lambda=\alpha+\beta\ne0$ in $\F_p$ and let $\ell$ be the
multiplicative order of a root of $t^p-t+\lambda$ in $\F_{p^p}$.
The coefficients of $F_{r,s,1}$ modulo $p^2$, indexed from $n=1$,
have minimal period
\[
 \begin{cases}
 \ell,&\beta=0\text{ in }\F_p,\\
 p\ell,&\beta\ne0\text{ in }\F_p.
 \end{cases}
\]
In particular, $\ell$ divides $p^p-1$.  It need not be assumed equal
to $p^p-1$; that equality in the A396807 case was proved by the
certificates above.
\end{proposition}
\begin{proof}
The Artin--Schreier factor is irreducible of degree $p$ by the
Frobenius-orbit calculation following Theorem \ref{thm:firstcorrection}.
Cancel $1-x$ from $1-x^{p-1}$ in \eqref{eq:Hgeneral}.  The remaining
numerator factors have degrees less than $p$, so none cancels the
irreducible factor.  At $x=1$ their value is a nonzero multiple of
$\beta$.  When $\beta\ne0$, the reduced denominator retains
$(1-x)^2$.  When $\beta=0$, $\alpha\ne0$ and one further factor $1-x$
cancels, leaving just a simple factor $1-x$.

The degree of the reduced numerator is at most that of its denominator.
Lemma \ref{lem:periodcriterion} now requires $\ell\mid L$ and, only in
the double-factor case, $p\mid L$.  Since $\ell\mid p^p-1$, it is prime
to $p$.  Formula \eqref{eq:firstlift} transfers exactly these periods
from $H$ over $\F_p$ to $F$ modulo $p^2$.
\end{proof}

\section{Exact compositional order and p-adic time}\label{sec:orders}

Periodicity in the coefficient index $n$ must not be confused with
periodicity in the iteration index $k$.  The latter has a different,
and particularly simple, exact answer.

\begin{lemma}[Raising a near-identity map]\label{lem:nearid}
If $G(x)\equiv x\pmod {p^e}$ with $e\geq1$, then
$\iter{G}{p}(x)\equiv x\pmod {p^{e+1}}$.
\end{lemma}
\begin{proof}
Write $G=x+p^e h$.  For any formal series $h$ with zero constant term,
$h(x+p^e h(x))-h(x)$ has all coefficients divisible by $p^e$.
Consequently composition of two near-identity maps adds their
$p^e$-perturbations modulo $p^{2e}$.  Since $2e\geq e+1$, induction gives
\[
 \iter{G}{j}\equiv x+jp^e h\pmod {p^{e+1}}.
\]
Taking $j=p$ proves the claim.  The argument includes $p=2$ and $e=1$.
\end{proof}

\begin{theorem}[Exact order]\label{thm:order}
For $F=F_{r,s,1}$ and any prime $p\mid D_{r,s}$, the order of the
reduction of $F$ in the composition group of
$x+x^2(\Z/p^e\Z)[[x]]$ is exactly $p^e$.  In particular,
\[
 \iter{A}{10^e}(x)\equiv x\pmod {10^e},
\]
and the compositional order of A396807 modulo $10^e$ is exactly $10^e$.
\end{theorem}
\begin{proof}
Modulo $p$, $F$ is $M$, which has $p$th iterate $x$.
Lemma \ref{lem:nearid} inductively gives
$\iter{F}{p^e}\equiv x\pmod {p^e}$.
The coefficient of $x^2$ in $F$ is $1$, and the coefficient of $x^2$
in its $k$th iterate is $k$.  Thus an iterate equal to $x$ modulo $p^e$
must have $p^e\mid k$, proving exactness.
For A396807 combine the two prime-power orders by the Chinese
remainder theorem, whose least common multiple is $10^e$.
\end{proof}

\begin{corollary}[A well-defined p-adic iteration parameter]\label{cor:padictime}
For each selected prime $p$, there is an injective group homomorphism
\[
 (\Z_p,+)\longrightarrow (x+x^2\Z_p[[x]],\circ),
 \qquad t\longmapsto\iter{F}{t},
\]
which extends the usual integer iterates, and
$\coeff{2}\iter{F}{t}=t$.
\end{corollary}
\begin{proof}
Choose integers $k_e$ representing $t$ modulo $p^e$.
Theorem \ref{thm:order} makes the reductions of $\iter{F}{k_e}$
compatible as $e$ increases, independently of the representatives.
Take the inverse limit coefficient by coefficient.  Every coefficient
of a composition is a finite polynomial in input coefficients, so the
integer identity
$\iter{F}{a}\circ\iter{F}{b}=\iter{F}{a+b}$ passes to the inverse limit.
The degree-two coefficient is $t$, proving injectivity.
\end{proof}

These are $p$-adic formal iterates, not real fractional iterates.
No assertion about a real-time flow follows from the construction.

\section{Two neighboring OEIS entries}\label{sec:neighbors}

\subsection{A396797: the same proof modulo 6}
The entry A396797 uses the equation
$B=x+\iter{B}{3}\iter{B}{4}$ and records both the constant coefficient
congruence and its all-integer-iterate version modulo $6$
\cite{oeis797}.  Here $D_{3,4}=6$, so Theorem \ref{thm:sharp} and
Corollary \ref{cor:iterates} immediately prove
\[
 B\equiv\frac{x}{1-x}\pmod6,
 \qquad \iter{B}{k}\equiv\frac{x}{1-kx}\pmod6
 \quad(k\in\Z).
\]
Thus this is not a second unrelated computation: it is another
specialization of the same exact-modulus theorem.

\subsection{A396798: a corrected start and all its modulo-8 iterates}
Let $G=F_{4,5,1}$, the generating function in A396798.  The entry lists
a repeating coefficient word modulo $8$ and several patterns for its
first eight iterates \cite{oeis798}.  A single formula handles all of
them, including a small indexing correction.

\begin{theorem}\label{thm:798}
Set $J(x)=x^4/(1-x)^3$.  For every integer $k$,
\begin{equation}\label{eq:798all}
 \iter{G}{k}(x)\equiv\frac{x}{1-kx}
       +4(k\bmod2)J(x)\pmod8.
\end{equation}
In particular,
\[
 G\equiv\frac{x}{1-x}+\frac{4x^4}{(1-x)^3}\pmod8.
\]
The coefficients of $G$ are $1$ at $n=1$, followed by the repeating
word $1,1,5,5$ starting at $n=2$.  For odd $k$, the coefficients of
$\iter{G}{k}$ starting at $n=2$ repeat
\begin{equation}\label{eq:oddword}
 k,\ 1,\ k+4,\ 5\pmod8.
\end{equation}
For even $k$, they are exactly the coefficients of $x/(1-kx)$ modulo $8$.
\end{theorem}
\begin{proof}
The exact uniform modulus is $D_{4,5}=4$, so $G\equiv M\pmod4$.
The residual of $M$, divided by $4$ and reduced modulo $2$, is
\[
 \frac{T(M)-M}{4}
 =\frac{2x^3-5x^4}{(1-x)(1-4x)(1-5x)}
 \equiv\frac{x^4}{(1-x)^2}\pmod2.
\]
For $(r,s,p)=(4,5,2)$, both $\alpha=2$ and $\beta=2$ vanish in $\F_2$,
so $C=0$ and \eqref{eq:L} reduces to $Lh=(1-x)h$.
The required correction is therefore $J=x^4/(1-x)^3$,
proving $G\equiv M+4J\pmod8$.

To iterate this identity, only the first variation modulo $2$ matters:
terms quadratic in $4J$ vanish modulo $8$.
Normalize $J=M^2w$, where
\[
 w=\frac{x^2}{1-x}.
\]
A direct substitution in $\F_2(x)$ gives $w(M)=w(x)$.
The sum in \eqref{eq:variation} is thus $(k\bmod2)w$.
For odd $k$, $M_k=M$ in $\F_2(x)$; for even $k$, the sum is zero.
This proves \eqref{eq:798all} for nonnegative $k$.
Theorem \ref{thm:order} gives $\iter{G}{8}=x$ modulo $8$, and both sides
of \eqref{eq:798all} depend only on $k$ modulo $8$, extending the formula
to all integers.

For $n\geq4$,
\[
 \coeff{n}J=\binom{n-2}{2};
\]
in smaller positive degrees it is zero.  Its parity, starting at
$n=2$, is the repeating word $0,0,1,1$.
For odd $k$, $k^{n-1}$ modulo $8$ alternates $k,1$ starting at $n=2$.
Adding the correction gives \eqref{eq:oddword}.
\end{proof}

For instance, the even iterates are
\[
 \iter{G}{2}\equiv x+2x^2+4x^3,\qquad
 \iter{G}{4}\equiv x+4x^2,\qquad
 \iter{G}{6}\equiv x+6x^2+4x^3,\qquad
 \iter{G}{8}\equiv x\pmod8.
\]
The odd words for $k=3,5,7$ are respectively
$(3,1,7,5)$, $(5,1,1,5)$, and $(7,1,3,5)$.

The literal proposed word for $G$ cannot start at $n=1$:
its third coefficient is $9\equiv1\pmod8$, not $5$.
The corrected start $n=2$ gives an exact theorem.  This is an indexing
correction in the neighboring entry, not a counterexample to either
of the two main A396807 conjectures.

\section{Large integer coefficients and zero complex radius}\label{sec:growth}

The rational reductions do not reflect convergence of the original
ordinary generating function.  In fact its coefficients grow too
quickly for any positive complex radius.

\begin{theorem}[A factorial lower bound]\label{thm:growth}
Let $F=F_{r,s,1}=\sum_{n\geq1}a_nx^n$ with $\max(r,s)\geq2$.
Then all coefficients are nonnegative, and
\begin{equation}\label{eq:2stepbound}
 a_n\geq(n-2)a_{n-2}\qquad(n\geq3).
\end{equation}
In particular, for $m\geq1$,
\begin{equation}\label{eq:factorialbounds}
 a_{2m}\geq2^{m-1}(m-1)!,\qquad
 a_{2m+1}\geq(2m-1)!!.
\end{equation}
The complex radius of convergence is zero.
\end{theorem}
\begin{proof}
The construction in Theorem \ref{thm:unique}, starting at $x$, uses only
products, compositions, and addition with nonnegative coefficients.
Hence the stabilized coefficients are nonnegative.  Also $F\geq x$
coefficientwise and $a_1=a_2=1$.  Coefficientwise order is preserved by
composition of such series, so any iterate of order at least two is
at least $F\circ F$, and every other iterate is at least $x$.
The defining equation thus implies
\[
 a_n\geq\coeff{n-1}(F\circ F).
\]
Within $F(F(x))$, select the outer term $a_{n-2}F(x)^{n-2}$.
Because $F=x+x^2+O(x^3)$, the coefficient of $x^{n-1}$ in
$F(x)^{n-2}$ is $n-2$: exactly one of its $n-2$ factors contributes
$x^2$ and all others contribute $x$.  This proves
\eqref{eq:2stepbound}.

Iterating it along the two parities gives \eqref{eq:factorialbounds}.
These lower bounds have $n$th roots tending to infinity.
For example, at least half the factors in $(m-1)!$ are at least
$(m-1)/2$, which already makes the $2m$th root unbounded.
The odd case is similar.  Therefore
$\limsup_n |a_n|^{1/n}=\infty$, and the radius is zero.
\end{proof}

This applies in particular to A396807.  A rational function regular at
zero has positive complex radius of convergence, so $A$ is not rational
over $\Q$.  There is no contradiction with rationality modulo every
$10^e$: reduction discards unbounded amounts of integer information,
and the rational representative is allowed to change with $e$.
The lower bound is not an asymptotic equivalent.  It also does not,
by itself, decide whether the series is D-finite.

\section{Algorithms, independent checks, and reproducibility}\label{sec:computation}

\subsection{An exact triangular construction}
Repeatedly recomputing the whole functional equation is unnecessary.
The supplied \texttt{series\_tools.py} maintains two tables:
\[
 P_{k,n}=\coeff{n}F(x)^k,
 \qquad B_{j,n}=\coeff{n}\iter{F}{j}(x).
\]
Here $P$ uses ordinary powers and $B$ uses functional iterates.
Initially $P_{0,0}=1$ and $B_{0,n}$ is $1$ at $n=1$ and $0$ otherwise.
At degree $n$, perform the following operations in order:
\begin{align}
 a_n&=\mathbf1_{n=1}
       +c\sum_{i=1}^{n-1}B_{r,i}B_{s,n-i},\label{eq:alg1}\\
 P_{1,n}&=a_n,\qquad
 P_{k,n}=\sum_{i=1}^{n-k+1}a_iP_{k-1,n-i}\quad(2\leq k\leq n),
                                                        \label{eq:alg2}\\
 B_{1,n}&=a_n,\qquad
 B_{j,n}=\sum_{k=1}^{n}B_{j-1,k}P_{k,n}
                    \quad(2\leq j\leq\max(r,s,1)).       \label{eq:alg3}
\end{align}
Equation \eqref{eq:alg1} uses only previously completed degrees because
both product factors have zero constant term.
Equation \eqref{eq:alg2} also uses only smaller degrees, and
\eqref{eq:alg3} is triangular in $j$.
The last identity represents $\iter{F}{j-1}\circ F$, which equals the
$j$th iterate.  These observations prove the algorithm directly.

Up to degree $N$, it uses
$O(N^3+\max(r,s)N^2)$ ring operations and
$O(N^2+\max(r,s)N)$ ring elements of storage.
These are arithmetic-operation bounds, not bit-complexity bounds;
integer sizes grow substantially.  Reducing after each operation gives
the same algorithm over any desired finite modulus.

For the last two digits alone, \eqref{eq:mod100} is much faster.
Expanding a fixed rational function gives a fixed-order linear
recurrence, hence $O(N)$ operations on bounded residues.
After storing one period, any particular index is answered by reducing
$n-1$ modulo $46860$ and adding one.  This does not require constructing
the huge integer $a(n)$.

\subsection{What was actually checked}
The standard-library verification program was run with Python 3.13.5.
The symbolic certificates were independently run with SymPy 1.14.0.
The generated files contain the full machine-readable logs.

\begin{center}
\begin{tabular}{@{}>{\raggedright\arraybackslash}p{0.56\textwidth}>{\raggedright\arraybackslash}p{0.35\textwidth}@{}}
\toprule
Check & Executed range / result\\
\midrule
Exact A396807 construction & $1\leq n\leq400$; $a(400)$ has 910 digits\\[3pt]
Published initial values & First 19 values matched\\[3pt]
Independent Horner-composition residual & Defining equation checked
 through degree 100\\[3pt]
Main and refined congruences & 2,000 coefficient checks through degree 400\\[3pt]
Positive and negative iterate congruences & $-20\leq k\leq20$,
 $1\leq n\leq48$: 1,968 checks\\[3pt]
Weighted sharp modulus and scaling & 1,014 parameter triples;
 29,406 coefficients\\[3pt]
Compositional order upper bounds & Moduli $2,4,8,16,5,25,125,625$
 through degree 40\\[3pt]
Neighboring entries & 100 A396797 checks; 580 A396798 checks\\[3pt]
Period data & 93,720 residues from the proved rational formulas;
 two full matching periods\\[3pt]
Symbolic identities & Universal residual, both first corrections,
 direct chain-rule checks, trace checks, and the A396798 correction passed\\
\bottomrule
\end{tabular}
\end{center}

The weighted family range was
$0\leq r,s\leq12$, $c\in\{-3,-1,0,1,2,5\}$, through degree $30$.
For every triple, the gcd of the computed discrepancies equaled
\eqref{eq:gcdideal}, and the scaling identity held coefficientwise.

The 400 large integer coefficients were independently generated; only
the first 19 were compared with the published initial list.
The 93,720 two-digit residues were obtained from the proved rational
formulas, \emph{not} by independently generating 93,720 full integer
coefficients.  This distinction prevents circular evidence: the
all-degree rational identities and period proof are established
algebraically above; the long residue file is an efficient consequence
and an implementation check.

As additional implementation checks, the program rejects every
maximal proper divisor of the proposed two-digit period.  It finds
the following discrepancies, all already at $n=1$:
\begin{center}
\begin{tabular}{@{}rrr@{}}
\toprule
Proposed shift $L$ & $a(1)\bmod100$ & $a(1+L)\bmod100$\\
\midrule
23430&01&81\\
15620&01&51\\
9372&01&21\\
4260&01&61\\
660&01&81\\
\bottomrule
\end{tabular}
\end{center}
These finite discrepancies are consistent with, but are not needed in
place of, the denominator and finite-field proof of minimality.

\subsection{Files and commands}
The archive contains the LaTeX source and compiled PDF, plus:
\begin{center}
\begin{tabular}{@{}p{0.44\textwidth}p{0.47\textwidth}@{}}
\toprule
File & Purpose\\
\midrule
\texttt{code/series\_tools.py} & Exact family solver, independent
 composition and inverse routines, rational recurrence expansion\\[3pt]
\texttt{code/verify.py} & Standard-library verification and data generation\\[3pt]
\texttt{code/symbolic\_checks.py} & Optional SymPy rational identities and
 reusable first-correction/linearized-solver functions\\[3pt]
\texttt{data/a396807\_exact\_400.txt} & Independently generated exact values\\[3pt]
\texttt{data/a396807\_mod100\_period.txt} & One complete minimal two-digit
 period, with indices\\[3pt]
\texttt{data/verification.json} & Detailed finite-test results\\[3pt]
\texttt{data/symbolic\_verification.json} & Symbolic certificate results\\[3pt]
\texttt{source\_notes/status.md} & Source provenance and status distinctions\\
\bottomrule
\end{tabular}
\end{center}
From the unpacked directory, run
\begin{lstlisting}
python code/verify.py
python code/symbolic_checks.py    # optional; requires SymPy
latexmk -pdf -interaction=nonstopmode article.tex
\end{lstlisting}
The first command requires no third-party Python packages.  The general
higher-power rational lifting procedure is given explicitly in the
proof of Theorem \ref{thm:rationallifts}; the code implements its
linearized inverse and closed first correction, not an optimized
arbitrary-exponent rational-function engine.

\section{Conclusions and further questions}\label{sec:conclusion}

The selected A396807 conjectures are proved, including negative
functional iterates and sharpness of the modulus.  The useful broader
lesson is that reduction to a finite-order rational substitution can
turn an equation involving unknown functional iterates into an
explicitly invertible trace equation.

For this family the chain of implications is fully proved:
\[
 \begin{aligned}
 \text{sharp M\"obius congruence}
 &\ \Longrightarrow\ \text{cyclic-trace inverse}\\
 &\ \Longrightarrow\ \text{rational prime-power lifts}\\
 &\ \Longrightarrow\ \text{eventually periodic residues}.
 \end{aligned}
\]
At the first additional digit the denominator is explicit enough to
prove exact periods.  For A396807 this produces a period of $46860$
for the last two digits, despite superexponential lower bounds on the
integer coefficients themselves.

Several questions remain beyond the results proved here.  One is to
obtain comparably explicit minimal periods modulo $10^e$ for $e\geq3$,
rather than merely eventual periodicity.  Another is to control the
growth of numerator and denominator degrees in the rational lifting
algorithm.  A third is to sharpen \eqref{eq:factorialbounds} to an actual
asymptotic description of $a(n)$.  Finally, the method as stated selects
primes dividing $D_{r,s}$; it does not classify reductions at other
primes.  These are natural continuation questions, not claims here
about their global status in the literature.

The computational evidence is deliberately separated from the proofs.
The finite checks test implementations and catch algebraic mistakes;
existence, congruence, rationality at every selected prime power, and
minimal period are established by the preceding arguments.

\begin{thebibliography}{9}
\bibitem{oeis807}
Paul D. Hanna, \emph{A396807: G.f. satisfies
$A(x)=x+A^5(x)A^6(x)$}, The On-Line Encyclopedia of Integer Sequences,
entry contributed June 16, 2026. Retrieved September 19, 2026.
\url{https://oeis.org/A396807}.

\bibitem{oeis807internal}
The OEIS Foundation, \emph{A396807, internal-format record},
record identifier \texttt{\#9 Jun 18 2026 00:43:54} in the retrieved
version. Retrieved September 19, 2026.
\url{https://oeis.org/A396807/internal}.

\bibitem{oeis797}
Paul D. Hanna, \emph{A396797: G.f. satisfies
$A(x)=x+A^3(x)A^4(x)$}, The On-Line Encyclopedia of Integer Sequences,
entry contributed June 15, 2026. Retrieved September 19, 2026.
\url{https://oeis.org/A396797}.

\bibitem{oeis798}
Paul D. Hanna, \emph{A396798: G.f. satisfies
$A(x)=x+A^4(x)A^5(x)$}, The On-Line Encyclopedia of Integer Sequences,
entry contributed June 16, 2026. Retrieved September 19, 2026.
\url{https://oeis.org/A396798}.
\end{thebibliography}

\end{document}
